% Part: computability
% Chapter: computability-theory
% Section: reducibility

\documentclass[../../../include/open-logic-section]{subfiles}

\begin{document}

\olfileid{cmp}{thy}{red}
\olsection{ન્યૂનીકરણક્ષમતા}

\begin{explain}
હવે આપણે જાણીએ છીએ કે ઓછામાં ઓછો એક ગણ $K_0$ એવો છે જે
સંગણકીય રીતે પરિગણનીય છે, પરંતુ સંગણનીય નથી. બીજા ગણો પણ
એવા જ હશે તે સ્પષ્ટ છે. ન્યૂનીકરણની પદ્ધતિ વારંવાર મૂળ
સિદ્ધાંતો તરફ પાછા જવું ન પડે તેમ બીજા ગણોમાં આ ગુણધર્મો
બતાવવાની શક્તિશાળી રીત આપે છે.

સામાન્ય રીતે, ગણ $A$નું ગણ~$B$માં ``ન્યૂનીકરણ'' એટલે
!!{element}s~$B$માં છે કે નહીં તે અંગેના જવાબોમાંથી
!!{element}s~$A$માં છે કે નહીં તે અંગેના જવાબો મેળવવાની
રીત. આપણે ``બહુ-એક ન્યૂનીકરણક્ષમતા'' પર ધ્યાન આપીશું,
જોકે જુદા ગુણધર્મો ધરાવતી ન્યૂનીકરણક્ષમતાની બીજી ઘણી
કલ્પનાઓ છે. સંગણકીય જટિલતાના અભ્યાસમાં પણ ન્યૂનીકરણની
કલ્પનાઓ કેન્દ્રસ્થાને છે; ત્યાં કાર્યક્ષમતાનો વિચાર પણ
કરવો પડે છે. દાખલા તરીકે, કોઈ ગણ NPમાં હોય અને દરેક NP
સમસ્યાને તેમાં એવી રીતથી ઘટાડી શકાય કે ન્યૂનીકરણ બહુ-એક
ન્યૂનીકરણ જેવું હોય, પરંતુ બહુપદી સમયમાં ગણી શકાય તેવી
વધારાની શરત પણ પૂરી કરે, તો તેને ``NP-પૂર્ણ'' કહે છે.

ન્યૂનીકરણની કલ્પના આપણે ગર્ભિત રીતે વાપરી જ છે. ગણ~$K$
આ રીતે વ્યાખ્યાયિત કરો:
\[
K = \Setabs{x}{\cfind{x}(x) \fdefined},
\]
એટલે $K = \Setabs{x}{x \in W_x}$. અટકવાની સમસ્યા
ઉકેલી શકાતી નથી તેની આપણી સાબિતી
(\olref[hlt]{thm:halting-problem}) સૌથી સીધે બતાવે છે કે
$K$ સંગણનીય નથી. યાદ કરો કે $K_0$ આ ગણ છે:
\[
K_0 = \Setabs{\tuple{e, x}}{\cfind{e}(x) \fdefined },
\]
એટલે $K_0 = \Setabs{\tuple{e,x}}{x \in W_e}$.
\sourcecorrection{OLCT-008}{સ્થિર અંગ્રેજી મૂળ બીજી
લખાવટમાં જોડીનો ક્રમ $\tuple{x,e}$ કરે છે. ઉપરની વ્યાખ્યા
અને $x \in W_e$ની શરત મુજબ ક્રમ $\tuple{e,x}$ જ છે.}
$K$ની અસંગણનીયતાની કોઈ પણ સાબિતીને $K_0$ની
અસંગણનીયતાની સાબિતી સુધી વિસ્તારી શકાય છે: જો $K_0$
સંગણનીય હોય, તો કોઈ !!a{element}~$x$, $K$માં છે કે નહીં
તે જોડી $\tuple{x, x}$, $K_0$માં છે કે નહીં એ પૂછીને નક્કી
કરી શકીએ. $x$ને $\tuple{x, x}$માં મોકલતું વિધેય~$f$,
$K$ના $K_0$માં \emph{ન્યૂનીકરણ}નું ઉદાહરણ છે.
\end{explain}

\begin{defn}
માનો કે $A$ અને $B$ પ્રાકૃતિક સંખ્યાઓના ગણો છે.
સંગણનીય વિધેય~$f\colon \Nat \to \Nat$, $A$નું~$B$માં
\emph{બહુ-એક ન્યૂનીકરણ} ત્યારે અને માત્ર ત્યારે છે જ્યારે
દરેક પ્રાકૃતિક સંખ્યા~$x$ માટે
\[
x \in A \quad \text{ત્યારે અને માત્ર ત્યારે} \quad f(x) \in B.
\]
આવું ન્યૂનીકરણ $f$ હોય, તો $A$ને~$B$માં \emph{બહુ-એક
ન્યૂનીકરણક્ષમ} કહીએ; તેને $A \leq_m B$ લખીએ. જો $A$
ના $B$માં અને વિરુદ્ધ દિશામાં પણ બહુ-એક ન્યૂનીકરણ હોય,
તો $A$ અને $B$ને \emph{બહુ-એક સમકક્ષ} કહીએ; તેને
$A \equiv_m B$ લખીએ.
\end{defn}

ઉપરની વ્યાખ્યાનું વિધેય $f$ એક-એક પણ હોય, તો $A$ને~$B$માં
\emph{એક-એક ન્યૂનીકરણક્ષમ} કહેવાય છે. નીચેના મોટા ભાગના
ન્યૂનીકરણો આ વધુ મજબૂત શરત પૂરી કરે છે, પરંતુ આ હકીકતનો
અહીં ઉપયોગ કરીશું નહીં.

\begin{digress}
એક-એક ન્યૂનીકરણક્ષમતા ખરેખર બહુ-એક ન્યૂનીકરણક્ષમતા કરતાં
વધુ મજબૂત શરત છે; આ સાચું છે, પણ બિલકુલ સ્પષ્ટ નથી.
બીજા શબ્દોમાં, એવા અનંત ગણો $A$ અને~$B$ છે કે $A$ને~$B$માં
બહુ-એક રીતે ઘટાડી શકાય છે, પરંતુ $B$માં એક-એક રીતે નહીં.
\end{digress}

\end{document}
